% Generated draft: requires mathematical and source review.
\documentclass[11pt]{article}
\usepackage[margin=25mm]{geometry}
\usepackage{fontspec}
\setmainfont{Latin Modern Roman}
\usepackage{amsmath,amssymb,mathtools,unicode-math}
\setmathfont{Latin Modern Math}
\usepackage{enumitem,needspace,xurl,hyperref}
\setlength{\parindent}{0pt}
\setlength{\parskip}{5pt}
\setlength{\emergencystretch}{3em}
\newcommand{\N}{\mathbb{N}}
\newcommand{\Z}{\mathbb{Z}}
\newcommand{\Q}{\mathbb{Q}}
\newcommand{\R}{\mathbb{R}}
\newcommand{\C}{\mathbb{C}}
\newcommand{\F}{\mathbb{F}}
\newcommand{\A}{\mathbb{A}}
\newcommand{\PP}{\mathbb{P}}
\newcommand{\E}{\mathbb{E}}
\newcommand{\eps}{\varepsilon}
\newcommand{\dd}{\,\mathrm{d}}
\begin{document}
\section*{UnsolvedMath ID: 30000156}
\textbf{Source problem code:} catalogue entry
\subsection{Definitions and mathematical statement}
\textbf{Conjecture 1, as supplied.} For each fixed birational map \(L\), let \(D_p(x)\) be the probability that a uniformly chosen point of \(\F_p^2\) belongs to a cycle of length at most \(px\). Does\[D(x)=\lim_{\substack{p\to\infty\\p\text{ prime}}}D_p(x)\] exist? The source attributes this conjecture to extensive experimental evidence, citing [6, 5, 7]. It does not specify the domain of \(x\), good-reduction convention, indeterminacy treatment, or whether existence is required at every \(x\) or only at continuity points of a limiting distribution. The same \(L\), defined over a suitable arithmetic coefficient field with specified reductions, must be used as \(p\) varies; unrelated maps selected separately for each prime do not define this question.\par Under a literal affine periodic-point convention, let \(L_p\) be the reduction at a permitted prime and let \(\operatorname{Per}_p\) consist of points of \(\F_p^2\) whose forward iterates are all defined and return to the starting point. For \(z\in\operatorname{Per}_p\), let \(\ell_p(z)\) be its least positive period. Then for \(x\ge0\),\[D_p(x)=\frac1{p^2}\#\{z\in\operatorname{Per}_p:\ell_p(z)\le\lfloor px\rfloor\}.\] This counts points on cycles, not cycles with equal weight, and normalizes by all \(p^2\) affine points, not just periodic points. For a polynomial permutation every point is periodic, but a general birational map can have indeterminacy and need not induce a permutation of \(\F_p^2\). A completion on a projective or resolved surface, a different exceptional-point convention, or conditioning on periodic points would require an explicitly different definition. The displayed convention is conditional and needs source confirmation.\par A pointwise limit at the cycle scale \(p\) need not by its existence alone be a probability distribution with total limiting mass one: long periods can escape that scale, and discontinuity points require special care. The source asks for existence, not universality of the same limiting function for all maps, and no tightness or rate is supplied.\par The supplied triage checked 2026-08-21 labels the problem partially solved. It reports that Roberts and Vivaldi's random-involution model explains limiting cycle statistics for important reversible polynomial maps, while explicitly cautioning that it does not prove the stated existence claim for every birational map. No theorem about the model is silently transferred to every deterministic reduction. This dated assessment requires independent verification.
\subsection{Short English statement}
For one fixed birational map reduced modulo growing primes, ask whether the fraction of affine points on cycles of length at most \(px\) has a limit. The map's arithmetic model and exceptional-point conventions are missing. Random-involution results for special reversible maps do not establish the universal claim.
\subsection{Sources}
\begin{itemize}\item Local UnsolvedMath registry, permanent ID 30000156, code catalogue entry; \emph{Limiting Cycle Distributions of Birational Maps}, Oberwolfach Reports Open Problems. Original source: Conjecture 1, with experimental references [6, 5, 7], whose bibliographic entries are absent.\item \emph{Finite Fields: Theory and Applications (2005)}, Oberwolfach Reports, DOI 10.4171/owr/2004/54, pp. 2913--2970. \url{https://doi.org/10.4171/owr/2004/54}. Proposed year and year tag: 2004; no proposer named.\item Supplied literature label \emph{Roberts--Vivaldi, 2009}: \url{https://arxiv.org/abs/0905.4135}.\item Supplied literature label \emph{original OWR report}: \url{https://doi.org/10.4171/owr/2004/54}.\item Literature triage checked 2026-08-21, status \texttt{partially\_solved}; record updated 2026-09-16. The \texttt{corrected\_verified} field concerns a concise restatement and balanced notation. Verify dated claims independently.\end{itemize}
\end{document}
